% =================================================================
% 3.3 Manage Moderation
% Feature description written against the real implementation:
%   apps/api/src/routes/forums/forums-admin/* and forums-reports/*
%   apps/web/app/admin/{moderation,reports}/*
% Procedures are described from the user interface point of view.
% =================================================================

\subsection{Moderation}

\subsubsection{View Moderation Queue}

\functionfield{Function Trigger} The logged-in administrator opens the Moderation screen by choosing \textbf{Forum} $\rightarrow$ \textbf{Moderation} in the admin sidebar.

\functionfield{Function description} This function allows an administrator to see all forum posts that are waiting to be reviewed. Every row shows the identifier of the post, its title, the automatic moderation score, the categories that were detected, the author and the creation date. The list only contains posts that are still pending, so the administrator always works on the items that need a decision.

\screenfield[0.95]{image/screenshot/sprint3/kvy/ViewModerationQueue.png}{Moderation queue screen}

\functionfield{Function details}
\begin{itemize}[leftmargin=1.6em,topsep=4pt,itemsep=3pt,parsep=0pt,label=$\bullet$]
    \item \detailhead{Precondition}
        \begin{itemize}[leftmargin=2.2em,topsep=2pt,itemsep=3pt,parsep=0pt,label=$\circ$]
            \item The user is logged in with an administrator account.
            \item The user is connected to the internet.
            \item The user is on the admin area and is allowed to open the Moderation screen.
        \end{itemize}
    \item \detailhead{Process}
        \begin{itemize}[leftmargin=2.2em,topsep=2pt,itemsep=3pt,parsep=0pt,label=$\circ$]
            \item \textbf{Open the Moderation Screen}: The user selects \textbf{Forum} in the sidebar and then \textbf{Moderation}, and the screen titled ``Moderation'' is displayed.
            \item \textbf{Read the Summary Line}: The subtitle shows the total number of posts that are waiting, in the form ``Review posts pending AI moderation (n pending)''.
            \item \textbf{Read Each Row}: Every row of the table presents the following information.
                \begin{itemize}[leftmargin=2.2em,topsep=2pt,itemsep=3pt,parsep=0pt,label=$\blacksquare$]
                    \item \textbf{ID}: The identifier of the post, shown as a small badge.
                    \item \textbf{Title}: The title of the post, shortened when it is too long.
                    \item \textbf{Score}: The automatic moderation score of the post, shown with two decimals, or ``N/A'' when no score was computed.
                    \item \textbf{Categories}: The categories detected by the automatic check, shown as badges, or ``---'' when none were detected.
                    \item \textbf{Author}: The display name of the user who wrote the post, or ``User \#id'' when the user has no display name.
                    \item \textbf{Created}: The date on which the post was created.
                \end{itemize}
            \item \textbf{Read the Score Colour}: The score badge uses a colour that reflects the risk level.
                \begin{itemize}[leftmargin=2.2em,topsep=2pt,itemsep=3pt,parsep=0pt,label=$\blacksquare$]
                    \item A score of 0.80 or more is shown in red, a score from 0.30 to 0.79 is shown in amber and a score below 0.30 is shown in green.
                \end{itemize}
            \item \textbf{Move Between Pages}: When there is more than one page, the user uses the page numbers or the \textbf{Previous} and \textbf{Next} buttons at the bottom of the list, and the counter shows ``Page x of y (n total)''.
            \item \textbf{Choose an Item}: The user opens, approves or rejects a post with the action buttons at the end of the row (see \textbf{View Moderation Item Detail} and \textbf{Resolve Moderation Item}).
        \end{itemize}
    \item \detailhead{Alternative Flow}
        \begin{itemize}[leftmargin=2.2em,topsep=2pt,itemsep=3pt,parsep=0pt,label=$\circ$]
            \item \textbf{No Pending Post}: When no post is waiting, the table shows the message ``No posts awaiting moderation.'' instead of rows.
            \item \textbf{Only One Page}: When all posts fit on a single page, the page numbers are hidden and only the table is shown.
        \end{itemize}
    \item \detailhead{Postconditions}
        \begin{itemize}[leftmargin=2.2em,topsep=2pt,itemsep=3pt,parsep=0pt,label=$\circ$]
            \item The administrator sees the complete list of posts that are waiting for a decision.
            \item No data is changed by this function, so the queue stays exactly as it was.
        \end{itemize}
    \item \detailhead{Error Handling}
        \begin{itemize}[leftmargin=2.2em,topsep=2pt,itemsep=3pt,parsep=0pt,label=$\circ$]
            \item \textbf{Not Allowed}: When a user who is not an administrator opens the screen, the user is sent to the access-denied page, because only administrators may read the moderation queue.
            \item \textbf{Lost Connection}: When the queue cannot be loaded, the screen shows an empty list so that the administrator can reload the page and try again.
        \end{itemize}
    \item \detailhead{Security Measures}
        \begin{itemize}[leftmargin=2.2em,topsep=2pt,itemsep=3pt,parsep=0pt,label=$\circ$]
            \item The whole admin area is protected: a guest or a banned account is sent to the login page, and a logged-in user without the administrator role is sent to the access-denied page.
            \item The moderation queue is only returned to administrator accounts, so ordinary users can never see posts that are waiting for review, no matter how they call the screen.
        \end{itemize}
\end{itemize}

% =================================================================

\subsubsection{View Moderation Item Detail}

\functionfield{Function Trigger} The logged-in administrator clicks the view button (the eye icon) at the end of a row in the moderation queue.

\functionfield{Function description} This function allows an administrator to read a pending post in full before making a decision. The dialog shows the title, the author, the creation date, the moderation score and the complete content of the post, together with the \textbf{Approve} and \textbf{Reject} buttons so that the decision can be taken from the same screen.

\screenfield[0.95]{image/screenshot/sprint3/kvy/ViewModerationItemDetail.png}{Moderation item detail dialog}

\functionfield{Function details}
\begin{itemize}[leftmargin=1.6em,topsep=4pt,itemsep=3pt,parsep=0pt,label=$\bullet$]
    \item \detailhead{Precondition}
        \begin{itemize}[leftmargin=2.2em,topsep=2pt,itemsep=3pt,parsep=0pt,label=$\circ$]
            \item The user is logged in with an administrator account.
            \item The user is connected to the internet.
            \item The user is on the Moderation screen and there is at least one pending post.
        \end{itemize}
    \item \detailhead{Process}
        \begin{itemize}[leftmargin=2.2em,topsep=2pt,itemsep=3pt,parsep=0pt,label=$\circ$]
            \item \textbf{Open the Detail}: The user clicks the view button on a row, and a dialog opens with the title of the post as its heading.
            \item \textbf{Read the Item Summary}: Under the title, the dialog shows the author of the post, the creation date and the moderation score (with ``N/A'' when no score exists).
            \item \textbf{Read the Content}: The dialog shows the whole content of the post with its formatting, exactly as it would appear to readers.
            \item \textbf{Scroll When Needed}: When the post is long, the dialog becomes scrollable so that the whole content can be read without leaving the screen.
            \item \textbf{Decide}: The user gives the decision with the \textbf{Approve} or \textbf{Reject} button at the bottom of the dialog (see \textbf{Resolve Moderation Item}).
            \item \textbf{Close the Detail}: The user closes the dialog with the close button or by clicking outside it, and returns to the queue.
        \end{itemize}
    \item \detailhead{Alternative Flow}
        \begin{itemize}[leftmargin=2.2em,topsep=2pt,itemsep=3pt,parsep=0pt,label=$\circ$]
            \item \textbf{No Score}: When the automatic check produced no score, the badge shows ``N/A'' and the rest of the dialog is unchanged.
            \item \textbf{Close Without Deciding}: When the user closes the dialog without pressing a button, no decision is recorded and the post stays in the queue.
        \end{itemize}
    \item \detailhead{Postconditions}
        \begin{itemize}[leftmargin=2.2em,topsep=2pt,itemsep=3pt,parsep=0pt,label=$\circ$]
            \item The administrator has read the complete post and knows its author, date, score and content.
            \item Nothing has changed while only viewing the item.
        \end{itemize}
    \item \detailhead{Error Handling}
        \begin{itemize}[leftmargin=2.2em,topsep=2pt,itemsep=3pt,parsep=0pt,label=$\circ$]
            \item \textbf{Not Allowed}: When a user who is not an administrator tries to open the detail, the action is refused because only administrators may review pending posts.
            \item \textbf{Item Already Handled}: When the post has already been approved or rejected by another administrator, the decision fails and a message tells the user that the post cannot be found or updated any more.
            \item \textbf{Lost Connection}: When the connection to the server is lost, the decision cannot be saved and the screen asks the user to try again.
        \end{itemize}
    \item \detailhead{Security Measures}
        \begin{itemize}[leftmargin=2.2em,topsep=2pt,itemsep=3pt,parsep=0pt,label=$\circ$]
            \item The detail is only available inside the admin area, so a pending post can never be read by an ordinary user.
            \item Every decision is recorded together with the administrator who made it and the moment it was made, so each review can be traced.
        \end{itemize}
\end{itemize}

% =================================================================

\subsubsection{Resolve Moderation Item}

\functionfield{Function Trigger} The logged-in administrator clicks the approve button or the reject button on a row of the moderation queue, or presses \textbf{Approve} or \textbf{Reject} inside the detail dialog.

\functionfield{Function description} This function allows an administrator to close a pending review by approving or rejecting a post. An approved post becomes visible to everyone in the forum. A rejected post stays hidden from the public and is moved to the trash, where it can be restored later. Before a rejection is applied, the administrator is asked to confirm it.

\screenfieldpair{image/screenshot/sprint3/kvy/ResolveModerationItem1.png}{image/screenshot/sprint3/kvy/ResolveModerationItem2.png}{Approve action (left) and reject confirmation (right)}

\functionfield{Function details}
\begin{itemize}[leftmargin=1.6em,topsep=4pt,itemsep=3pt,parsep=0pt,label=$\bullet$]
    \item \detailhead{Precondition}
        \begin{itemize}[leftmargin=2.2em,topsep=2pt,itemsep=3pt,parsep=0pt,label=$\circ$]
            \item The user is logged in with an administrator account.
            \item The user is connected to the internet.
            \item The post that is being resolved is still waiting in the moderation queue.
        \end{itemize}
    \item \detailhead{Process}
        \begin{itemize}[leftmargin=2.2em,topsep=2pt,itemsep=3pt,parsep=0pt,label=$\circ$]
            \item \textbf{Approve a Post}: The user clicks the approve button on a row, and the post is accepted straight away.
                \begin{itemize}[leftmargin=2.2em,topsep=2pt,itemsep=3pt,parsep=0pt,label=$\blacksquare$]
                    \item The post leaves the moderation queue and becomes visible to everyone in the forum.
                    \item The screen shows the message ``Post approved and now visible publicly''.
                \end{itemize}
            \item \textbf{Start a Rejection}: The user clicks the reject button on a row, or presses \textbf{Reject} inside the detail dialog.
            \item \textbf{Confirm the Rejection}: When the reject button of a row is used, a confirmation dialog titled ``Reject Post'' appears with the question ``Reject ``title''? The post will not be visible publicly and will be moved to trash.''.
                \begin{itemize}[leftmargin=2.2em,topsep=2pt,itemsep=3pt,parsep=0pt,label=$\blacksquare$]
                    \item The user presses \textbf{Reject} to confirm, or \textbf{Cancel} to keep the post in the queue.
                \end{itemize}
            \item \textbf{Apply the Rejection}: When the rejection is confirmed, the post is refused and is moved to the trash.
                \begin{itemize}[leftmargin=2.2em,topsep=2pt,itemsep=3pt,parsep=0pt,label=$\blacksquare$]
                    \item The post is never shown to ordinary users while it stays rejected.
                    \item The screen shows the message ``Post rejected and moved to trash''.
                \end{itemize}
            \item \textbf{Update the Queue}: The resolved post disappears from the list and the remaining items move up, so the queue always shows only the posts that still need a decision.
        \end{itemize}
    \item \detailhead{Alternative Flow}
        \begin{itemize}[leftmargin=2.2em,topsep=2pt,itemsep=3pt,parsep=0pt,label=$\circ$]
            \item \textbf{Cancel the Rejection}: When the user closes the confirmation dialog, the post stays pending and nothing changes.
            \item \textbf{Reject from the Detail Dialog}: When \textbf{Reject} is pressed inside the detail dialog, the rejection is applied immediately without an extra confirmation dialog.
            \item \textbf{Empty Queue}: When the last pending post is resolved, the table shows the message ``No posts awaiting moderation.''.
        \end{itemize}
    \item \detailhead{Postconditions}
        \begin{itemize}[leftmargin=2.2em,topsep=2pt,itemsep=3pt,parsep=0pt,label=$\circ$]
            \item The post is no longer pending and has been marked as approved or rejected.
            \item An approved post is visible in the forum, while a rejected post is hidden and can be found in the trash.
            \item The decision is recorded with the administrator who made it and the moment it was made.
        \end{itemize}
    \item \detailhead{Error Handling}
        \begin{itemize}[leftmargin=2.2em,topsep=2pt,itemsep=3pt,parsep=0pt,label=$\circ$]
            \item \textbf{Not Allowed}: When a user who is not an administrator tries to approve or reject a post, the action is refused because only administrators may resolve moderation items.
            \item \textbf{Item Not Found}: When the post has already been resolved by another administrator, the action fails and the screen shows the message ``Couldn't approve post. Try again.'' or ``Couldn't reject post. Try again.''.
            \item \textbf{Cannot Save the Decision}: When the decision cannot be saved, the post stays in the queue and a message asks the administrator to try again.
            \item \textbf{Lost Connection}: When the connection to the server is lost, the decision is not applied and the screen asks the user to try again.
        \end{itemize}
    \item \detailhead{Security Measures}
        \begin{itemize}[leftmargin=2.2em,topsep=2pt,itemsep=3pt,parsep=0pt,label=$\circ$]
            \item Only administrators may approve or reject a post, and the check is made again on the server for every decision.
            \item The post can only be resolved while it is still pending, so a decision cannot be applied twice to the same post.
            \item A rejected post is hidden from the public and kept in the trash instead of being erased, so the decision can be reviewed and undone later.
        \end{itemize}
\end{itemize}


% =================================================================
% 3.4 Manage ID Verification
% Feature description written against the real implementation:
%   apps/api/src/routes/identity-verifications/identity-verifications-admin/*
%   apps/web/app/admin/identity-verifications/*
% Procedures are described from the user interface point of view.
% =================================================================

\subsection{ID Verification}

\subsubsection{View Verification Request List}

\functionfield{Function Trigger} The logged-in administrator opens the ID Verification screen by choosing \textbf{Users} $\rightarrow$ \textbf{Identity Verification} in the admin sidebar.

\functionfield{Function description} This function allows an administrator to see all the identity verification requests that users have submitted. Every row shows the identifier, the user, the type of document, the status and the submission date. The administrator can filter the requests by status and search by email or display name.

\screenfield[0.95]{image/screenshot/sprint3/kvy/ViewVerificationRequestList.png}{ID verification request list screen}

\functionfield{Function details}
\begin{itemize}[leftmargin=1.6em,topsep=4pt,itemsep=3pt,parsep=0pt,label=$\bullet$]
    \item \detailhead{Precondition}
        \begin{itemize}[leftmargin=2.2em,topsep=2pt,itemsep=3pt,parsep=0pt,label=$\circ$]
            \item The user is logged in with an administrator account.
            \item The user is connected to the internet.
            \item At least one user has already submitted an identity document, otherwise the list is empty.
        \end{itemize}
    \item \detailhead{Process}
        \begin{itemize}[leftmargin=2.2em,topsep=2pt,itemsep=3pt,parsep=0pt,label=$\circ$]
            \item \textbf{Open the ID Verification Screen}: The user selects \textbf{Users} in the sidebar and then \textbf{Identity Verification}, and the screen titled ``ID Verification Requests'' is displayed.
            \item \textbf{Read the Summary Line}: The subtitle shows the total number of requests in the form ``Review and approve user identity documents (n total)''.
            \item \textbf{Read Each Row}: Every row of the table presents the following information.
                \begin{itemize}[leftmargin=2.2em,topsep=2pt,itemsep=3pt,parsep=0pt,label=$\blacksquare$]
                    \item \textbf{ID}: The identifier of the request, shown as a small badge.
                    \item \textbf{Name}: The avatar, the display name and the email of the user who submitted the document.
                    \item \textbf{Document Type}: The kind of document, shown as National ID card, Passport, Driver license, Residence permit or Other document.
                    \item \textbf{Status}: The current status, shown as a badge with the labels Pending, Approved or Rejected.
                    \item \textbf{Submitted}: The date on which the document was submitted.
                \end{itemize}
            \item \textbf{Read the Status Colour}: A pending request is shown in amber, an approved request in green and a rejected request in red.
            \item \textbf{Filter by Status}: The user chooses one of the tabs \textbf{All}, \textbf{Pending}, \textbf{Approved} or \textbf{Rejected} to show only the requests with that status.
            \item \textbf{Search}: The user types a text in the search box with the placeholder ``Search by email or display name...'' and submits it, and only the matching requests are shown.
            \item \textbf{Move Between Pages}: When there is more than one page, the user uses the page numbers or the \textbf{Previous} and \textbf{Next} buttons at the bottom of the list.
            \item \textbf{Choose a Request}: The user opens the request with the view button (see \textbf{View ID Verification Request Detail}).
        \end{itemize}
    \item \detailhead{Alternative Flow}
        \begin{itemize}[leftmargin=2.2em,topsep=2pt,itemsep=3pt,parsep=0pt,label=$\circ$]
            \item \textbf{No Request}: When no request matches the current filter or search, the table shows the message ``No verification requests found.''.
            \item \textbf{Clear the Search}: When the user clears the search box, all requests are listed again.
        \end{itemize}
    \item \detailhead{Postconditions}
        \begin{itemize}[leftmargin=2.2em,topsep=2pt,itemsep=3pt,parsep=0pt,label=$\circ$]
            \item The administrator sees the requests that match the chosen status and search text.
            \item No data is changed by this function, so the requests stay exactly as they were.
        \end{itemize}
    \item \detailhead{Error Handling}
        \begin{itemize}[leftmargin=2.2em,topsep=2pt,itemsep=3pt,parsep=0pt,label=$\circ$]
            \item \textbf{Not Allowed}: When a user who is not an administrator opens the screen, the user is sent to the access-denied page, because only administrators may read identity documents.
            \item \textbf{Lost Connection}: When the requests cannot be loaded, the screen shows an empty list so that the administrator can reload the page and try again.
        \end{itemize}
    \item \detailhead{Security Measures}
        \begin{itemize}[leftmargin=2.2em,topsep=2pt,itemsep=3pt,parsep=0pt,label=$\circ$]
            \item The list of requests is only returned to administrator accounts, so the identity data of a user is never exposed to other users.
            \item Only the basic information is listed, while the images of the document are shown only inside the detail of a request.
        \end{itemize}
\end{itemize}

% =================================================================

\subsubsection{View ID Verification Request Detail}

\functionfield{Function Trigger} The logged-in administrator clicks the view button (the eye icon) at the end of a row in the ID verification request list.

\functionfield{Function description} This function allows an administrator to examine an identity verification request in full. The dialog shows the user, the submission date, the type of document, the front and back images of the document, the review notes if a decision has already been made, and the buttons that let the administrator approve or reject a pending request.

\screenfield[0.95]{image/screenshot/sprint3/kvy/ViewIDVerificationRequestDetail.png}{ID verification request detail dialog}

\functionfield{Function details}
\begin{itemize}[leftmargin=1.6em,topsep=4pt,itemsep=3pt,parsep=0pt,label=$\bullet$]
    \item \detailhead{Precondition}
        \begin{itemize}[leftmargin=2.2em,topsep=2pt,itemsep=3pt,parsep=0pt,label=$\circ$]
            \item The user is logged in with an administrator account.
            \item The user is connected to the internet.
            \item The user is on the ID Verification screen and the request exists.
        \end{itemize}
    \item \detailhead{Process}
        \begin{itemize}[leftmargin=2.2em,topsep=2pt,itemsep=3pt,parsep=0pt,label=$\circ$]
            \item \textbf{Open the Detail}: The user clicks the view button on a row, and a dialog titled ``ID Verification Request \#id'' opens.
            \item \textbf{Read the Header}: The header shows the avatar, the display name and the email of the user, together with the badge of the document type and the badge of the status.
            \item \textbf{Read the Submission Date}: The line ``Submitted \dots'' shows the moment at which the document was sent, including the time.
            \item \textbf{Read the Document}: The section \textbf{Document} shows the front side of the document and, when it was uploaded, the back side as well.
                \begin{itemize}[leftmargin=2.2em,topsep=2pt,itemsep=3pt,parsep=0pt,label=$\blacksquare$]
                    \item When the back side was not submitted, a dashed placeholder shows the message ``No back side uploaded with this submission.''.
                \end{itemize}
            \item \textbf{Zoom an Image}: The user clicks an image to open it in full size in a larger dialog, and closes the larger dialog to come back to the request.
            \item \textbf{Read the Review Notes}: When a decision has already been made, the section \textbf{Review Notes} shows the note that was left by the administrator who handled the request.
            \item \textbf{Decide on a Pending Request}: When the request is pending, the dialog also shows a review notes box and the \textbf{Approve} and \textbf{Reject} buttons (see \textbf{Approve or Reject ID Verification}).
            \item \textbf{Close the Detail}: The user closes the dialog and returns to the list, where the row reflects the current status.
        \end{itemize}
    \item \detailhead{Alternative Flow}
        \begin{itemize}[leftmargin=2.2em,topsep=2pt,itemsep=3pt,parsep=0pt,label=$\circ$]
            \item \textbf{Already Decided}: When the request has already been approved or rejected, the dialog only shows the images, the status and the review notes, and no decision button appears.
            \item \textbf{No Back Side}: When only the front side exists, the dialog shows the placeholder instead of the second image and the rest is unchanged.
            \item \textbf{Close Without Deciding}: When the user closes the dialog without pressing a button, no decision is recorded and the request stays pending.
        \end{itemize}
    \item \detailhead{Postconditions}
        \begin{itemize}[leftmargin=2.2em,topsep=2pt,itemsep=3pt,parsep=0pt,label=$\circ$]
            \item The administrator has seen the user, the document, the date and any previous review notes.
            \item Nothing has changed while only viewing the request.
        \end{itemize}
    \item \detailhead{Error Handling}
        \begin{itemize}[leftmargin=2.2em,topsep=2pt,itemsep=3pt,parsep=0pt,label=$\circ$]
            \item \textbf{Not Allowed}: When a user who is not an administrator tries to open the detail, the action is refused because only administrators may read identity documents.
            \item \textbf{Request Not Found}: When the request no longer exists, the detail cannot be opened and a message tells the user that the request cannot be found.
            \item \textbf{Document Not Available}: When an image of the document can no longer be read, the image is not shown and a message tells the user that the document is not available.
            \item \textbf{Lost Connection}: When the detail cannot be loaded, the screen asks the user to try again.
        \end{itemize}
    \item \detailhead{Security Measures}
        \begin{itemize}[leftmargin=2.2em,topsep=2pt,itemsep=3pt,parsep=0pt,label=$\circ$]
            \item The images of the document are delivered only to an administrator account, so the document of a user is never exposed to anyone else.
            \item The document images are opened inside the admin area as read-only content, so they cannot be saved or edited from the screen.
            \item The decision buttons are hidden as soon as the request has been handled, so a request can be reviewed only while it is pending.
        \end{itemize}
\end{itemize}

% =================================================================

\subsubsection{Approve or Reject ID Verification}

\functionfield{Function Trigger} The logged-in administrator presses the \textbf{Approve} or \textbf{Reject} button inside the detail of a pending ID verification request.

\functionfield{Function description} This function allows an administrator to close an identity verification request by approving or rejecting the document that the user submitted. The administrator can write an optional review note that is stored with the decision. Approving a request marks the user as verified, while rejecting a request leaves the user unverified.

\screenfield[0.8]{image/screenshot/sprint3/kvy/ApprovedOrRejectedIDVerification.png}{Approving or rejecting an ID verification request}

\functionfield{Function details}
\begin{itemize}[leftmargin=1.6em,topsep=4pt,itemsep=3pt,parsep=0pt,label=$\bullet$]
    \item \detailhead{Precondition}
        \begin{itemize}[leftmargin=2.2em,topsep=2pt,itemsep=3pt,parsep=0pt,label=$\circ$]
            \item The user is logged in with an administrator account.
            \item The user is connected to the internet.
            \item The request is pending and its document has been examined in the detail dialog.
        \end{itemize}
    \item \detailhead{Process}
        \begin{itemize}[leftmargin=2.2em,topsep=2pt,itemsep=3pt,parsep=0pt,label=$\circ$]
            \item \textbf{Write a Review Note}: The user may type a note in the box labeled ``Review Notes'', which is optional.
                \begin{itemize}[leftmargin=2.2em,topsep=2pt,itemsep=3pt,parsep=0pt,label=$\blacksquare$]
                    \item The note is stored with the decision and is shown again the next time the request is opened.
                \end{itemize}
            \item \textbf{Approve the Request}: The user presses \textbf{Approve}, and the button shows ``Approving\dots'' while the decision is being saved.
                \begin{itemize}[leftmargin=2.2em,topsep=2pt,itemsep=3pt,parsep=0pt,label=$\blacksquare$]
                    \item The request is marked as approved, the user is marked as verified and the screen shows the message ``Identity verification approved''.
                \end{itemize}
            \item \textbf{Reject the Request}: The user presses \textbf{Reject}, and the button shows ``Rejecting\dots'' while the decision is being saved.
                \begin{itemize}[leftmargin=2.2em,topsep=2pt,itemsep=3pt,parsep=0pt,label=$\blacksquare$]
                    \item The request is marked as rejected, the user stays unverified and the screen shows the message ``Identity verification rejected''.
                \end{itemize}
            \item \textbf{See the Result}: When the decision is saved, the status badge of the request changes and the decision buttons disappear, so the request is now read-only.
        \end{itemize}
    \item \detailhead{Alternative Flow}
        \begin{itemize}[leftmargin=2.2em,topsep=2pt,itemsep=3pt,parsep=0pt,label=$\circ$]
            \item \textbf{No Review Note}: When the user leaves the box empty, the decision is saved without a note.
            \item \textbf{Note Too Long}: When the note is longer than the allowed length, the note is refused and the administrator is asked to shorten it before the decision can be saved.
            \item \textbf{Reopen the Detail Later}: When the request is opened again, it shows the chosen status, the review note and no decision buttons.
        \end{itemize}
    \item \detailhead{Postconditions}
        \begin{itemize}[leftmargin=2.2em,topsep=2pt,itemsep=3pt,parsep=0pt,label=$\circ$]
            \item The request has the status that the administrator chose, and it is no longer pending.
            \item An approved user is marked as verified, while a rejected user stays unverified.
            \item The decision, the note and the moment of the decision are stored with the request.
        \end{itemize}
    \item \detailhead{Error Handling}
        \begin{itemize}[leftmargin=2.2em,topsep=2pt,itemsep=3pt,parsep=0pt,label=$\circ$]
            \item \textbf{Not Allowed}: When a user who is not an administrator tries to approve or reject a request, the action is refused because only administrators may decide on identity documents.
            \item \textbf{Request Not Found}: When the request no longer exists, the decision is refused and the screen shows the message ``Couldn't update the request. Try again.''.
            \item \textbf{Already Reviewed}: When another administrator has already handled the request, the decision is refused and a message tells the user that the request has already been reviewed.
            \item \textbf{Lost Connection}: When the connection to the server is lost, the decision is not saved and the screen asks the user to try again.
        \end{itemize}
    \item \detailhead{Security Measures}
        \begin{itemize}[leftmargin=2.2em,topsep=2pt,itemsep=3pt,parsep=0pt,label=$\circ$]
            \item Only administrators may approve or reject a request, and the check is made again on the server for every decision.
            \item A request can only be decided while it is still pending, so a decision cannot be applied twice to the same request.
            \item Each decision is recorded together with the administrator who made it, the note and the moment of the decision, so every review can be traced.
        \end{itemize}
\end{itemize}

% =================================================================
% 3.5 Manage Forum Trash
% Feature description written against the real implementation:
%   apps/api/src/routes/forums/forums-admin/*
%   apps/web/app/admin/trash/*
% Procedures are described from the user interface point of view.
% =================================================================

\subsection{Trash}

\subsubsection{View Deleted Post List}

\functionfield{Function Trigger} The logged-in administrator opens the Trash screen by choosing \textbf{Forum} $\rightarrow$ \textbf{Trash} in the admin sidebar while the \textbf{Posts} tab is selected.

\functionfield{Function description} This function allows an administrator to see all the posts that have been removed from the forum, including the posts that were rejected during moderation. Every row shows the identifier, the title, the status, the author and the moment of deletion, so the administrator can decide which posts to restore.

\screenfield[0.95]{image/screenshot/sprint3/kvy/ViewDeletedPostList.png}{Deleted post list screen}

\functionfield{Function details}
\begin{itemize}[leftmargin=1.6em,topsep=4pt,itemsep=3pt,parsep=0pt,label=$\bullet$]
    \item \detailhead{Precondition}
        \begin{itemize}[leftmargin=2.2em,topsep=2pt,itemsep=3pt,parsep=0pt,label=$\circ$]
            \item The user is logged in with an administrator account.
            \item The user is connected to the internet.
            \item At least one post has been deleted or rejected, otherwise the trash is empty.
        \end{itemize}
    \item \detailhead{Process}
        \begin{itemize}[leftmargin=2.2em,topsep=2pt,itemsep=3pt,parsep=0pt,label=$\circ$]
            \item \textbf{Open the Trash Screen}: The user selects \textbf{Forum} in the sidebar and then \textbf{Trash}, and the screen titled ``Trash'' is displayed with the tab \textbf{Posts} active.
            \item \textbf{Read the Summary Line}: The subtitle explains that the trash contains deleted posts and comments, including the posts that were rejected during moderation.
            \item \textbf{Read Each Row}: Every row of the table presents the following information.
                \begin{itemize}[leftmargin=2.2em,topsep=2pt,itemsep=3pt,parsep=0pt,label=$\blacksquare$]
                    \item \textbf{ID}: The identifier of the post, shown as a small badge.
                    \item \textbf{Title}: The title of the post, shortened when it is too long.
                    \item \textbf{Status}: The moderation status of the post, shown as the badge Rejected, Pending or Approved.
                    \item \textbf{Author}: The avatar and the display name of the user who wrote the post, or ``User \#id'' when the user has no display name.
                    \item \textbf{Deleted on}: The date on which the post was deleted, or ``---'' when the post is in the trash because it was rejected.
                \end{itemize}
            \item \textbf{Switch to Comments}: The user presses the \textbf{Comments} tab at the top right of the screen to see the deleted comments instead (see \textbf{View Deleted Comment List}).
            \item \textbf{Move Between Pages}: When there is more than one page, the user uses the page numbers or the \textbf{Previous} and \textbf{Next} buttons at the bottom of the list.
            \item \textbf{Choose a Post}: The user opens a post with the view button, or restores it with the restore button in the row (see \textbf{View Deleted Post Detail} and \textbf{Restore Deleted Post}).
        \end{itemize}
    \item \detailhead{Alternative Flow}
        \begin{itemize}[leftmargin=2.2em,topsep=2pt,itemsep=3pt,parsep=0pt,label=$\circ$]
            \item \textbf{Empty Trash}: When nothing has been deleted or rejected, the table shows the message ``Trash is empty.''.
            \item \textbf{Rejected Post Without a Deletion Date}: When a post is in the trash because it was rejected, the column shows ``---'' while the status badge shows Rejected.
        \end{itemize}
    \item \detailhead{Postconditions}
        \begin{itemize}[leftmargin=2.2em,topsep=2pt,itemsep=3pt,parsep=0pt,label=$\circ$]
            \item The administrator sees all the posts that are currently out of the forum.
            \item No data is changed by this function, so the trash stays exactly as it was.
        \end{itemize}
    \item \detailhead{Error Handling}
        \begin{itemize}[leftmargin=2.2em,topsep=2pt,itemsep=3pt,parsep=0pt,label=$\circ$]
            \item \textbf{Not Allowed}: When a user who is not an administrator opens the screen, the user is sent to the access-denied page, because only administrators may read the trash.
            \item \textbf{Lost Connection}: When the trash cannot be loaded, the screen shows an empty list so that the administrator can reload the page and try again.
        \end{itemize}
    \item \detailhead{Security Measures}
        \begin{itemize}[leftmargin=2.2em,topsep=2pt,itemsep=3pt,parsep=0pt,label=$\circ$]
            \item The trash is only returned to administrator accounts, so deleted and rejected posts are never visible to ordinary users.
            \item A rejected post is kept in the trash instead of being erased, so it can be restored later if the decision was wrong.
        \end{itemize}
\end{itemize}

% =================================================================

\subsubsection{View Deleted Post Detail}

\functionfield{Function Trigger} The logged-in administrator clicks the view button (the eye icon) at the end of a row in the deleted post list.

\functionfield{Function description} This function allows an administrator to read a deleted or rejected post in full before restoring it. The dialog shows the title, the status, the author, the creation date, the moment of deletion and the complete content of the post.

\screenfield[0.95]{image/screenshot/sprint3/kvy/ViewDeletedPostDetail.png}{Deleted post detail dialog}

\functionfield{Function details}
\begin{itemize}[leftmargin=1.6em,topsep=4pt,itemsep=3pt,parsep=0pt,label=$\bullet$]
    \item \detailhead{Precondition}
        \begin{itemize}[leftmargin=2.2em,topsep=2pt,itemsep=3pt,parsep=0pt,label=$\circ$]
            \item The user is logged in with an administrator account.
            \item The user is connected to the internet.
            \item The user is on the Trash screen with the \textbf{Posts} tab active and the post exists in the trash.
        \end{itemize}
    \item \detailhead{Process}
        \begin{itemize}[leftmargin=2.2em,topsep=2pt,itemsep=3pt,parsep=0pt,label=$\circ$]
            \item \textbf{Open the Detail}: The user clicks the view button on a row, and a dialog opens with the title of the post and its status badge as its heading.
            \item \textbf{Read the Summary}: Under the title, the dialog shows the author of the post, the creation date, the moment of deletion and the moderation score.
                \begin{itemize}[leftmargin=2.2em,topsep=2pt,itemsep=3pt,parsep=0pt,label=$\blacksquare$]
                    \item When the post is in the trash because it was rejected and not because it was deleted, the line reads ``In trash (rejected)'' instead of the deletion date.
                \end{itemize}
            \item \textbf{Read the Content}: The dialog shows the whole content of the post with its formatting, exactly as it was written by the author.
            \item \textbf{Scroll When Needed}: When the post is long, the dialog becomes scrollable so that the whole content can be read.
            \item \textbf{Close the Detail}: The user closes the dialog and returns to the trash list.
        \end{itemize}
    \item \detailhead{Alternative Flow}
        \begin{itemize}[leftmargin=2.2em,topsep=2pt,itemsep=3pt,parsep=0pt,label=$\circ$]
            \item \textbf{No Score}: When the automatic check produced no score, the badge shows ``N/A'' and the rest of the dialog is unchanged.
            \item \textbf{Rejected Post}: When the post was rejected during moderation, the dialog shows the Rejected badge and the line ``In trash (rejected)''.
        \end{itemize}
    \item \detailhead{Postconditions}
        \begin{itemize}[leftmargin=2.2em,topsep=2pt,itemsep=3pt,parsep=0pt,label=$\circ$]
            \item The administrator has read the complete post and knows why it is out of the forum.
            \item Nothing has changed while only viewing the post.
        \end{itemize}
    \item \detailhead{Error Handling}
        \begin{itemize}[leftmargin=2.2em,topsep=2pt,itemsep=3pt,parsep=0pt,label=$\circ$]
            \item \textbf{Not Allowed}: When a user who is not an administrator tries to open the detail, the action is refused because only administrators may read the trash.
            \item \textbf{Post Not Found}: When the post is no longer in the trash, the detail cannot be opened and a message tells the user that the post cannot be found.
            \item \textbf{Lost Connection}: When the detail cannot be loaded, the screen asks the user to try again.
        \end{itemize}
    \item \detailhead{Security Measures}
        \begin{itemize}[leftmargin=2.2em,topsep=2pt,itemsep=3pt,parsep=0pt,label=$\circ$]
            \item The detail is only available inside the admin area, so the content of a deleted post is never shown to ordinary users.
            \item The dialog is read-only, so the content of a deleted post cannot be edited from this screen.
        \end{itemize}
\end{itemize}

% =================================================================

\subsubsection{Restore Deleted Post}

\functionfield{Function Trigger} The logged-in administrator clicks the restore button (the round arrow icon) at the end of a row in the deleted post list.

\functionfield{Function description} This function allows an administrator to bring a post back into the forum. A post that was only deleted becomes visible again straight away, while a post that was rejected goes back to the moderation queue so that it can be reviewed again. Before the post is restored, the administrator is asked to confirm the action.

\screenfield[0.8]{image/screenshot/sprint3/kvy/RestoreDeletedPost.png}{Restore post confirmation dialog}

\functionfield{Function details}
\begin{itemize}[leftmargin=1.6em,topsep=4pt,itemsep=3pt,parsep=0pt,label=$\bullet$]
    \item \detailhead{Precondition}
        \begin{itemize}[leftmargin=2.2em,topsep=2pt,itemsep=3pt,parsep=0pt,label=$\circ$]
            \item The user is logged in with an administrator account.
            \item The user is connected to the internet.
            \item The post is currently in the trash, either because it was deleted or because it was rejected during moderation.
        \end{itemize}
    \item \detailhead{Process}
        \begin{itemize}[leftmargin=2.2em,topsep=2pt,itemsep=3pt,parsep=0pt,label=$\circ$]
            \item \textbf{Start the Restore}: The user clicks the restore button on a row, and a confirmation dialog titled ``Restore Post'' appears.
            \item \textbf{Read the Confirmation}: The dialog shows the question ``Restore ``title'' back to the forum?'' and explains what will happen.
                \begin{itemize}[leftmargin=2.2em,topsep=2pt,itemsep=3pt,parsep=0pt,label=$\blacksquare$]
                    \item When the post was rejected, the dialog adds ``It will return to the moderation queue for review.''.
                    \item When the post was only deleted, the dialog adds ``It will be visible publicly again immediately.''.
                \end{itemize}
            \item \textbf{Confirm the Restore}: The user presses \textbf{Restore}, and the button shows ``Restoring\dots'' while the action is being applied.
            \item \textbf{Apply the Restore}: The post is taken out of the trash and the screen shows the message ``Post restored successfully''.
                \begin{itemize}[leftmargin=2.2em,topsep=2pt,itemsep=3pt,parsep=0pt,label=$\blacksquare$]
                    \item A post that was only deleted becomes visible in the forum again straight away.
                    \item A post that was rejected is put back into the moderation queue with the status Pending, so the administrator can review it again in the Moderation screen.
                \end{itemize}
            \item \textbf{Update the Trash}: The restored post disappears from the trash list and the remaining rows move up.
        \end{itemize}
    \item \detailhead{Alternative Flow}
        \begin{itemize}[leftmargin=2.2em,topsep=2pt,itemsep=3pt,parsep=0pt,label=$\circ$]
            \item \textbf{Cancel the Restore}: When the user presses \textbf{Cancel} or closes the dialog, the post stays in the trash and nothing changes.
            \item \textbf{Empty Trash}: When the last post is restored, the table shows the message ``Trash is empty.''.
            \item \textbf{Restore Several Posts}: The user repeats the action on other rows, one post at a time.
        \end{itemize}
    \item \detailhead{Postconditions}
        \begin{itemize}[leftmargin=2.2em,topsep=2pt,itemsep=3pt,parsep=0pt,label=$\circ$]
            \item The post is no longer in the trash.
            \item A restored post that was only deleted is visible in the forum again.
            \item A restored post that was rejected is waiting again in the moderation queue with the status Pending and without a previous review.
        \end{itemize}
    \item \detailhead{Error Handling}
        \begin{itemize}[leftmargin=2.2em,topsep=2pt,itemsep=3pt,parsep=0pt,label=$\circ$]
            \item \textbf{Not Allowed}: When a user who is not an administrator tries to restore a post, the action is refused because only administrators may restore content.
            \item \textbf{Post Not Found}: When the post is no longer in the trash, the restore fails and the screen shows the message ``Couldn't restore post. Try again.''.
            \item \textbf{Cannot Restore}: When the post cannot be restored, it stays in the trash and a message asks the administrator to try again.
            \item \textbf{Lost Connection}: When the connection to the server is lost, the restore is not applied and the screen asks the user to try again.
        \end{itemize}
    \item \detailhead{Security Measures}
        \begin{itemize}[leftmargin=2.2em,topsep=2pt,itemsep=3pt,parsep=0pt,label=$\circ$]
            \item Only administrators may restore a post, and the check is made again on the server for every restore.
            \item Only a post that is really in the trash can be restored, so the action cannot be used on a post that is already visible in the forum.
            \item A restored post that had been rejected is never made public directly: it is put back into the moderation queue, so it must pass a new review before it can be seen by users.
        \end{itemize}
\end{itemize}

% =================================================================

\subsubsection{View Deleted Comment List}

\functionfield{Function Trigger} The logged-in administrator opens the Trash screen by choosing \textbf{Forum} $\rightarrow$ \textbf{Trash} in the admin sidebar and then presses the \textbf{Comments} tab.

\functionfield{Function description} This function allows an administrator to see all the comments that have been removed from the forum. Every row shows the identifier, the text of the comment, the post it belongs to, the author and the moment of deletion, so the administrator can decide which comments to restore.

\screenfield[0.95]{image/screenshot/sprint3/kvy/ViewDeletedCommentList.png}{Deleted comment list screen}

\functionfield{Function details}
\begin{itemize}[leftmargin=1.6em,topsep=4pt,itemsep=3pt,parsep=0pt,label=$\bullet$]
    \item \detailhead{Precondition}
        \begin{itemize}[leftmargin=2.2em,topsep=2pt,itemsep=3pt,parsep=0pt,label=$\circ$]
            \item The user is logged in with an administrator account.
            \item The user is connected to the internet.
            \item At least one comment has been deleted, otherwise the trash is empty.
        \end{itemize}
    \item \detailhead{Process}
        \begin{itemize}[leftmargin=2.2em,topsep=2pt,itemsep=3pt,parsep=0pt,label=$\circ$]
            \item \textbf{Open the Comments Tab}: The user opens the Trash screen and presses the \textbf{Comments} tab at the top right, and the screen titled ``Trash'' is displayed with the comments of the trash.
            \item \textbf{Read Each Row}: Every row of the table presents the following information.
                \begin{itemize}[leftmargin=2.2em,topsep=2pt,itemsep=3pt,parsep=0pt,label=$\blacksquare$]
                    \item \textbf{ID}: The identifier of the comment, shown as a small badge.
                    \item \textbf{Comment}: The text of the comment, shortened when it is too long.
                    \item \textbf{Post}: The title of the post that the comment belongs to, or ``Post \#id'' when the title is not available.
                    \item \textbf{Author}: The avatar and the display name of the user who wrote the comment, or ``User \#id'' when the user has no display name.
                    \item \textbf{Deleted on}: The date on which the comment was deleted.
                \end{itemize}
            \item \textbf{Switch to Posts}: The user presses the \textbf{Posts} tab at the top right of the screen to see the deleted posts instead (see \textbf{View Deleted Post List}).
            \item \textbf{Move Between Pages}: When there is more than one page, the user uses the page numbers or the \textbf{Previous} and \textbf{Next} buttons at the bottom of the list.
            \item \textbf{Choose a Comment}: The user opens a comment with the view button, or restores it with the restore button in the row (see \textbf{View Deleted Comment Detail} and \textbf{Restore Deleted Comment}).
        \end{itemize}
    \item \detailhead{Alternative Flow}
        \begin{itemize}[leftmargin=2.2em,topsep=2pt,itemsep=3pt,parsep=0pt,label=$\circ$]
            \item \textbf{Empty Trash}: When no comment has been deleted, the table shows the message ``Trash is empty.''.
            \item \textbf{Missing Post Title}: When the post of the comment is not available, the row shows only the identifier of the post.
        \end{itemize}
    \item \detailhead{Postconditions}
        \begin{itemize}[leftmargin=2.2em,topsep=2pt,itemsep=3pt,parsep=0pt,label=$\circ$]
            \item The administrator sees all the comments that are currently out of the forum.
            \item No data is changed by this function, so the trash stays exactly as it was.
        \end{itemize}
    \item \detailhead{Error Handling}
        \begin{itemize}[leftmargin=2.2em,topsep=2pt,itemsep=3pt,parsep=0pt,label=$\circ$]
            \item \textbf{Not Allowed}: When a user who is not an administrator opens the screen, the user is sent to the access-denied page, because only administrators may read the trash.
            \item \textbf{Lost Connection}: When the trash cannot be loaded, the screen shows an empty list so that the administrator can reload the page and try again.
        \end{itemize}
    \item \detailhead{Security Measures}
        \begin{itemize}[leftmargin=2.2em,topsep=2pt,itemsep=3pt,parsep=0pt,label=$\circ$]
            \item The list of deleted comments is only returned to administrator accounts, so a deleted comment is never shown to ordinary users.
            \item The delete action of a comment does not erase its content, so the administrator can still read what was removed before deciding to restore it.
        \end{itemize}
\end{itemize}

% =================================================================

\subsubsection{View Deleted Comment Detail}

\functionfield{Function Trigger} The logged-in administrator clicks the view button (the eye icon) at the end of a row in the deleted comment list.

\functionfield{Function description} This function allows an administrator to read a deleted comment in full before restoring it. The dialog shows the post that the comment belongs to, the author, the complete text of the comment, the identifier of the comment and the moment of deletion.

\screenfield[0.95]{image/screenshot/sprint3/kvy/ViewDeletedCommentDetail.png}{Deleted comment detail dialog}

\functionfield{Function details}
\begin{itemize}[leftmargin=1.6em,topsep=4pt,itemsep=3pt,parsep=0pt,label=$\bullet$]
    \item \detailhead{Precondition}
        \begin{itemize}[leftmargin=2.2em,topsep=2pt,itemsep=3pt,parsep=0pt,label=$\circ$]
            \item The user is logged in with an administrator account.
            \item The user is connected to the internet.
            \item The user is on the Trash screen with the \textbf{Comments} tab active and the comment exists in the trash.
        \end{itemize}
    \item \detailhead{Process}
        \begin{itemize}[leftmargin=2.2em,topsep=2pt,itemsep=3pt,parsep=0pt,label=$\circ$]
            \item \textbf{Open the Detail}: The user clicks the view button on a row, and a dialog titled ``Deleted Comment'' opens.
            \item \textbf{Read the Summary}: Under the title, the dialog shows the title of the post that the comment belongs to and the name of the author.
            \item \textbf{Read the Content}: The dialog shows the complete text of the comment, exactly as it was written by the author.
            \item \textbf{Read the Footer}: The footer of the dialog shows the identifier of the comment and the moment at which it was deleted.
            \item \textbf{Close the Detail}: The user closes the dialog and returns to the trash list.
        \end{itemize}
    \item \detailhead{Alternative Flow}
        \begin{itemize}[leftmargin=2.2em,topsep=2pt,itemsep=3pt,parsep=0pt,label=$\circ$]
            \item \textbf{Missing Post Title}: When the title of the post is not available, the dialog shows only the identifier of the post.
            \item \textbf{Missing Deletion Date}: When the deletion date is not available, the footer shows ``---''.
        \end{itemize}
    \item \detailhead{Postconditions}
        \begin{itemize}[leftmargin=2.2em,topsep=2pt,itemsep=3pt,parsep=0pt,label=$\circ$]
            \item The administrator has read the complete comment and knows to which post it belongs.
            \item Nothing has changed while only viewing the comment.
        \end{itemize}
    \item \detailhead{Error Handling}
        \begin{itemize}[leftmargin=2.2em,topsep=2pt,itemsep=3pt,parsep=0pt,label=$\circ$]
            \item \textbf{Not Allowed}: When a user who is not an administrator tries to open the detail, the action is refused because only administrators may read the trash.
            \item \textbf{Comment Not Found}: When the comment is no longer in the trash, the detail cannot be opened and a message tells the user that the comment cannot be found.
            \item \textbf{Lost Connection}: When the detail cannot be loaded, the screen asks the user to try again.
        \end{itemize}
    \item \detailhead{Security Measures}
        \begin{itemize}[leftmargin=2.2em,topsep=2pt,itemsep=3pt,parsep=0pt,label=$\circ$]
            \item The detail is only available inside the admin area, so the content of a deleted comment is never shown to ordinary users.
            \item The dialog is read-only, so the content of a deleted comment cannot be edited from this screen.
        \end{itemize}
\end{itemize}

% =================================================================

\subsubsection{Restore Deleted Comment}

\functionfield{Function Trigger} The logged-in administrator clicks the restore button (the round arrow icon) at the end of a row in the deleted comment list.

\functionfield{Function description} This function allows an administrator to bring a deleted comment back to its post. The comment is shown again to users when its post is still visible in the forum. Before the comment is restored, the administrator is asked to confirm the action.

\screenfield[0.8]{image/screenshot/sprint3/kvy/RestoreDeletedComment.png}{Restore comment confirmation dialog}

\functionfield{Function details}
\begin{itemize}[leftmargin=1.6em,topsep=4pt,itemsep=3pt,parsep=0pt,label=$\bullet$]
    \item \detailhead{Precondition}
        \begin{itemize}[leftmargin=2.2em,topsep=2pt,itemsep=3pt,parsep=0pt,label=$\circ$]
            \item The user is logged in with an administrator account.
            \item The user is connected to the internet.
            \item The comment is currently in the trash because it was deleted.
        \end{itemize}
    \item \detailhead{Process}
        \begin{itemize}[leftmargin=2.2em,topsep=2pt,itemsep=3pt,parsep=0pt,label=$\circ$]
            \item \textbf{Start the Restore}: The user clicks the restore button on a row, and a confirmation dialog titled ``Restore Comment'' appears.
            \item \textbf{Read the Confirmation}: The dialog shows the question ``Restore this comment on post ``title''? It will be visible again (as long as its post is not deleted).''.
            \item \textbf{Confirm the Restore}: The user presses \textbf{Restore}, and the button shows ``Restoring\dots'' while the action is being applied.
            \item \textbf{Apply the Restore}: The comment is taken out of the trash and the screen shows the message ``Comment restored successfully''.
                \begin{itemize}[leftmargin=2.2em,topsep=2pt,itemsep=3pt,parsep=0pt,label=$\blacksquare$]
                    \item The comment appears again under its post when the post itself is still visible in the forum.
                    \item When the post has also been deleted, the comment stays hidden until the post is restored as well.
                \end{itemize}
            \item \textbf{Update the Trash}: The restored comment disappears from the trash list and the remaining rows move up.
        \end{itemize}
    \item \detailhead{Alternative Flow}
        \begin{itemize}[leftmargin=2.2em,topsep=2pt,itemsep=3pt,parsep=0pt,label=$\circ$]
            \item \textbf{Cancel the Restore}: When the user presses \textbf{Cancel} or closes the dialog, the comment stays in the trash and nothing changes.
            \item \textbf{Deleted Post}: When the post of the comment has also been deleted, the comment is restored but remains hidden until the post comes back to the forum.
            \item \textbf{Empty Trash}: When the last comment is restored, the table shows the message ``Trash is empty.''.
        \end{itemize}
    \item \detailhead{Postconditions}
        \begin{itemize}[leftmargin=2.2em,topsep=2pt,itemsep=3pt,parsep=0pt,label=$\circ$]
            \item The comment is no longer in the trash and is visible again under its post when the post is visible.
            \item The rest of the comments and the post itself are unchanged.
        \end{itemize}
    \item \detailhead{Error Handling}
        \begin{itemize}[leftmargin=2.2em,topsep=2pt,itemsep=3pt,parsep=0pt,label=$\circ$]
            \item \textbf{Not Allowed}: When a user who is not an administrator tries to restore a comment, the action is refused because only administrators may restore content.
            \item \textbf{Comment Not Found}: When the comment is no longer in the trash, the restore fails and the screen shows the message ``Couldn't restore comment. Try again.''.
            \item \textbf{Cannot Restore}: When the comment cannot be restored, it stays in the trash and a message asks the administrator to try again.
            \item \textbf{Lost Connection}: When the connection to the server is lost, the restore is not applied and the screen asks the user to try again.
        \end{itemize}
    \item \detailhead{Security Measures}
        \begin{itemize}[leftmargin=2.2em,topsep=2pt,itemsep=3pt,parsep=0pt,label=$\circ$]
            \item Only administrators may restore a comment, and the check is made again on the server for every restore.
            \item Only a comment that is really in the trash can be restored, so the action cannot be used on a comment that is already visible.
            \item A restored comment can never become visible outside a visible post, so the owner of a deleted post cannot see comments that belong to removed content.
        \end{itemize}
\end{itemize}
